\subsection{Strong relational systems and clique atlases}
We prove Theorem~\ref{thm:strong-relations}, including the
partial-order assertions for infinite figures.

Start with the carrier presentation. Internal refinement is
reflexive and antisymmetric because refinement is a partial order.
If $c=c(x,y)$, then $x\rho c$: any smaller candidate below $c$
would also contradict its leastness below $y$.
If $x\rho z\leq y$, then $c(x,y)\leq z$, and the carrier
condition below $z$ forces $c(x,y)=z$. This proves uniqueness
of factorization. Applying the carrier construction to
$x'\leq x\rho y$ proves exchange.

It remains to justify internal transitivity. If $x\rho y\rho z$,
write $x\rho c\leq z$ for its carrier factorization.
If $c\ne z$, principal disjointness gives $c\delta z$.
Heredity along $y\rho z$ and $x\rho y$ gives $c\delta x$;
heredity along $x\rho c$ then gives $x\delta x$, a contradiction.
Thus $c=z$ and $x\rho z$.

Conversely start with the factorization presentation and define
$\ll$ as in the theorem. Reflexivity and containment of incidence
are immediate. For transitivity, from
$x\rho a\leq y\rho b\leq z$, exchange gives
$a\rho c\leq b$, and then $x\rho c\leq z$.
For antisymmetry suppose also $y\rho b\leq x$.
The same exchange gives $x\rho a\rho c\leq x$.
Uniqueness of factorization into $x$, compared with
$x\rho x\leq x$, gives $c=x$. Antisymmetry of $\rho$
gives $a=x$, hence $x\leq y$. Reversing the roles gives
$y\leq x$, so $x=y$.
If $x\rho c\leq y$ and $x\ll d\leq y$, factorization
through $d$ gives the same $c$, so $c\leq d$.
Thus $c$ is the least carrier and the recovered internal relation
is exactly the original one. The constructions are inverse.

Now let $\mathscr H$ be an admitted clique atlas as in the theorem.
Bounded unions are admitted by downward-closed subfigure closure
inside their bound. Every figure is the union of its principal
downsets. Those principals are completely join-prime by membership
of their maximum, and this union expression shows that they are
exactly the nonzero completely join-irreducible figures.
Two figures are compatible exactly when their union is a clique.
For principal downsets this says
\[
 x\bowtie y\quad\Longleftrightarrow\quad
 (\forall a\leq x,\ a\not\leq y)
 (\forall b\leq y,\ b\not\leq x)\ a\delta b.
\]
Pairs lying in one principal downset already satisfy the required
disjunction; these are precisely the remaining pairs.

A stratum $p$ refining a figure $F$ has a carrier: factor its
refinement to some $x\in F$ as $p\rho c\leq x$.
Two possible carriers $c,d\in F$, if distinct, would satisfy
$c\delta d$, and heredity would imply $p\delta p$.
They therefore agree and give the least refined stratum in $F$.
If $P,F$ are compatible and $P\ll F$, the carrier $c$ of
$p\in P$ lies in their common clique. If $p\ne c$,
$p\delta c$ and $p\rho c$ again give a self-disjunction.
Thus $p=c\in F$ and $P\subseteq F$.

Elementwise refinement is reflexive and transitive. Suppose
$P\ll F\ll P$. For $p\in P$, take its successive carriers
$q\in F,r\in P$. Then $p\rho q\rho r$.
The carrier of $p$ in $P$ is $p$, so internal transitivity and
uniqueness give $r=p$. Antisymmetry of $\rho$ yields $q=p$.
Hence $P\subseteq F$; symmetry gives equality. This proves
refinement antisymmetry without finiteness. The remaining
carrier-atelier clauses now follow directly from the definitions.

\begin{proposition}[The abstract restriction calculus]
\label{prop:abstract-carriers}
In any prime carrier atelier, carriers compose. If $P\ll F$
and $G\leq F$, the strata of $P$ whose $F$-carriers lie in
$\widetilde G$ form a figure $P|G$. It is the greatest common
lower bound of $P,G$ for refinement. Within one parent,
restriction commutes with incidence unions and intersections.
\end{proposition}
\begin{proof}
Prime generation identifies the incidence ideal below a figure
with all downward-closed subsets of its stratum deployment:
the join of such a subset has exactly that deployment by primality.
If $x\ll P\ll F$, put $y=c_P(x)$ and $z=c_F(y)$.
Then $x\rho y\rho z$, so uniqueness gives $z=c_F(x)$.
If $x'\leq x$, the carrier of $x$ is a candidate refined by
$x'$, whence $c_F(x')\leq c_F(x)$. Thus the specified inverse
image of $\widetilde G$ is downward closed and defines a figure.
It refines both $P$ and $G$.

If $Q\ll P,G$ and $q\in\widetilde Q$, let $p=c_P(q)$.
Composition gives $c_F(p)=c_F(q)$. Since $q\ll G\leq F$,
the latter carrier lies in $\widetilde G$.
Thus $p\in\widetilde{P|G}$ and $q\ll p$; the stratum
test gives $Q\ll P|G$. This proves the infimum property.
The union and intersection assertions are the corresponding
inverse-image identities for the carrier map.
\end{proof}

For the clique atlas, let $F,G$ be compatible parents,
$K=F\cap G$, and $P\ll F,Q\ll G$.
Compatibility of $P|K,Q|K$ implies compatibility of $P,Q$.
Indeed, for distinct $p\in P,q\in Q$, compare their parent
carriers $c,d$. If $c\ne d$, the common parent clique gives
$c\delta d$, hence $p\delta q$. If $c=d$, it belongs to
$K$, and the required disjunction is provided by the two
compatible restrictions.
If those restrictions are also disjoint as sets, then $P,Q$
are disjoint: a common stratum would have two parent carriers;
distinct ones force self-disjunction, while equal ones put it
in both restrictions. In particular disjoint parents have
disjoint refinements. Distinct compatible strata likewise have
pairwise disjoint open lieux, by heredity along their internal
refinements.

For the gluing statement of Theorem~\ref{thm:strong-relations},
the matching restrictions make every pair $P_i,P_j$ compatible
by the preceding argument. Their union $P$ is admitted for a
finite cover, and for any cover in the full atlas. It refines $F$.
If $p\in P_j$ has carrier in $F_i$, that carrier lies in
$F_i\cap F_j$; matching puts $p$ in $P_i$.
Therefore $P|F_i=P_i$. Conversely every stratum of any proposed
refinement of $F$ is carried by some covering $F_i$, so these
restrictions determine it. This proves uniqueness.
A directed incidence union is a clique downset because any two
of its strata lie together in one member, proving the final
full-atlas assertion.

\subsection{The representation by admissible set figures}
We prove Theorem~\ref{thm:set-atlas} in both directions.
Start with a prime carrier atelier and use its canonical sets
\[
 A_x=\{z:z\ll x\},\qquad \Phi_F=\{A_x:x\leq F\}.
\]
One has $A_x\subseteq A_y$ if and only if $x\ll y$;
thus the sets distinguish the strata. For a point $z$ in the
support of $\Phi_F$, its least containing member is
$A_{c_F(z)}$. The cell of $A_x$ is consequently
\[
 U_x=\{z:z\rho x\},
\]
independent of the parent figure and nonempty because $x\rho x$.
Hence $\Phi_F$ is admissible.
Inside one figure, inclusion $A_x\subseteq A_y$ is compatible
refinement and therefore means $x\leq y$. The principal
subfigure at $A_x$ is exactly $\Phi_x$.

Incidence is preserved and reflected by prime deployment.
If $F\ll G$ and $x\leq F$ has carrier $y$ in $G$,
internal transitivity gives $U_x\subseteq U_y$, proving
set-figure refinement. Conversely $U_x\subseteq U_y$
puts $x\in U_y$ and hence $x\rho y$; the stratum test
recovers the original refinement.

Prime generation turns bounded joins into unions, giving
condition~1. The principal identity gives condition~2.
Inclusion $A_x\subseteq A_y$ is refinement of $x$ to $y$,
so the preceding argument gives condition~3. Compatibility
detected on strata gives condition~4. Finally the global
family is all principal refinement ideals $A_x$. A point $x$
has least member $A_x$, and that member has global residual
$\{x\}$. These sets cover $M$ and separate points by
antisymmetry of refinement. This proves condition~5.

Conversely suppose the five set-atlas conditions hold.
Global admissibility gives each point $e$ a least member $A_e$.
It belongs to a global stratum $B$ precisely when $A_e\subseteq B$.
Two points with the same least member therefore have the same
memberships; separation makes them equal. Nonempty residuals
give a unique point $e_A$ for every occurring stratum $A$.
Thus points and occurring strata are in bijection.

Order $\mathscr H$ by inclusion. Bounded unions exist by
condition~1. By condition~2 every figure is the union of its
admitted principal subfigures $\widehat A$. Each is completely
join-prime: membership of its maximum $A$ in a union occurs
in one summand, which then contains all of $\widehat A$.
The union expression shows that these are exactly the nonzero
completely join-irreducibles. Condition~4 gives compatibility
detected on strata. A closed subfamily is admitted because it
is a bounded union of its principal subfigures.

Principal consistency makes the cell $C_A$ of $A$ independent
of its admitted parent. Set-figure refinement is reflexive and
transitive. If two figures refine each other, their nonempty
partition cells coincide. Equality of a cell of $A$ with one
of $B$ puts $e_A\in B$ and $e_B\in A$, giving both
$A\subseteq B$ and $B\subseteq A$. Thus $A=B$ and
the figures coincide. Refinement is a partial order. Inclusion
of figures preserves their cells and hence implies refinement.

If a cell of $A$ is contained in a cell of $B$, then
$e_A\in B$, so $A\subseteq B$. Condition~3 gives
$\widehat A\ll\widehat B$. Conversely, refinement of
$\widehat A$ to a principal subfigure of $\Psi$ maps its
top cell into a cell of $\Psi$. This proves exactly the
elementwise refinement axiom.

For $\widehat A\ll\Psi$, let $B$ be the least member of
$\Psi$ containing $e_A$. Global leastness gives $A\subseteq B$,
so condition~3 gives $\widehat A\ll\widehat B$.
If $D\in\Psi$ is another refined stratum, then $e_A\in D$,
and leastness in $\Psi$ gives $B\subseteq D$ and
$B\in\widehat D$. Hence $B$ is the least carrier.

Suppose compatible figures $\Phi,\Psi$ satisfy
$\Phi\ll\Psi$. In a common admitted parent their cells
are the same as their own cells. A nonempty cell of $\Phi$
contained in a cell of $\Psi$ must be that same parent cell.
Their strata agree, giving $\Phi\subseteq\Psi$.
Finally $A\rho B$ is equivalent to $C_A\subseteq C_B$:
the forward implication uses the carrier of the top cell,
and the reverse uses condition~3 and the least carrier just
proved. Cell containment is transitive, establishing the last
prime carrier axiom.

The two constructions recover the objects canonically.
On the abstract side they identify $x$ with $\Phi_x$ and
recover both orders. On the set side the bijection $e\mapsto A_e$
identifies the canonical refinement ideal of $B$ with its
original set, because $e_A\in B$ exactly when $A\subseteq B$.
The admitted figures are recovered as well. Atelier isomorphisms
transport strata and refinement ideals; bijections of set-atlases
transport cells and their containment. The two recovery maps
commute with these isomorphisms. This proves the stated groupoid
equivalence, including the empty ground set.

\begin{counterexample}[Why the principal subfigure must be admitted]
\label{ex:missing-principal}
On $E=\{a,b\}$ take
\[
 \Phi=\{\{a\},\{a,b\}\},\qquad
 \mathscr H=\{\varnothing,\Phi\}.
\]
This satisfies the five conditions of Theorem~\ref{thm:set-atlas}
with the membership requirement deleted from condition~2 and
all principal refinements interpreted in the ambient
$\operatorname{Fig}(E)$, but it does not recover the prescribed
two ground objects as strata.
\end{counterexample}
\begin{proof}
Bounded unions and consistency of shared principal subfigures
are immediate. The only strict inclusion of strata is
$\{a\}\subset\{a,b\}$, whose principal singleton figure
refines $\Phi$ in the ambient family. For the pair $\Phi,\Phi$,
condition~4 has false premise at the pair $\{a\},\{a\}$:
its principal singleton is not admitted. Pairs involving the
empty figure satisfy the conclusion. The global family is
admissible and separating, with cells $\{a\}$ and $\{b\}$.
Yet $\widehat{\{a\}}=\{\{a\}\}\notin\mathscr H$,
and the incidence poset has only one nonzero irreducible.
\end{proof}

Admissibility also needs an actual least stratum at every point
when infinite families are allowed. On
$E=\mathbb N\cup\{\infty\}$, let
$A_n=\{\infty\}\cup\{m:m\geq n\}$.
Every residual is the nonempty singleton $\{n\}$ and
$A_n\cap A_m=A_{\max(n,m)}$, but $\infty$ has no least
containing member. Thus nonempty residuals and binary-intersection
coverage alone are insufficient. If binary intersections are
unions of common lower strata, adding coverage by the residuals
is sufficient: a residual point of $A$ in another stratum $B$
lies in a common lower stratum; minimality forces that stratum
to be $A$, giving $A\subseteq B$. Every residual point then
has its required least member, and nonempty residuals ensure
that every stratum occurs.

\subsection{Local detection and the original compatibility}
We prove Theorem~\ref{thm:local-relations}. A lieu $p$ has
singleton principal deployment and $U_p=\{p\}$. Every lieu
refining a figure has a carrier in it, so
\begin{equation}\label{eq:lieu-deployment}
 \ell(F)=\bigcup_{x\leq F}U_x.
\end{equation}
For two admitted clique figures,
\begin{equation}\label{eq:figure-disjoint-pairs}
 F\parallel G\quad\Longleftrightarrow\quad
 (\forall x\in F)(\forall y\in G)\ x\delta y.
\end{equation}
The right side excludes a common stratum by irreflexivity and
makes the union an admitted clique. Conversely a compatible
disjoint union has every cross pair distinct and disjoint.
In particular $p\parallel q$ is exactly $p\delta q$ for lieux.

Start with the locally detected strong presentation. If $U_x$
were empty, local detection for $x,x$ would give $x\delta x$,
so condition~1 follows. Substituting local detection into
\eqref{eq:figure-disjoint-pairs}, and then using
\eqref{eq:lieu-deployment}, proves condition~2, even for
arbitrary admitted figures. Distinct compatible strata are
members of one clique; hereditary disjunction along their open
lieux gives condition~3.

For condition~4, let $a\leq x$ but $a\not\leq y$.
Every $p\in U_a$ lies outside $\ell(K)$: otherwise its carrier
$c$ in $K$ would be another internal carrier in $x$, forcing
$a=c\leq y$. Likewise $b\leq y$ but $b\not\leq x$ has
$U_b\subseteq\ell(y)\setminus\ell(K)$.
The hypothesis of condition~4 and local detection give
$a\delta b$. These are exactly the cross pairs not already
contained in one principal clique. Their union is therefore
admitted, proving compatibility of $x,y$.

Conversely retain the two orders and carriers of the atelier
and define $\delta$ by the open-lieu formula with graph
$p\parallel q$. Symmetry is immediate. Nonempty $U_x$ and
$p\not\parallel p$ make it irreflexive. Internal transitivity
gives $U_{x'}\subseteq U_x$ whenever $x'\rho x$, proving
heredity. Condition~3 gives all-pairs principal disjointness.
These are all the strong-system axioms.

We check that this construction recovers the original
compatibility. If $x,y$ were compatible, all distinct strata
in their common deployment satisfy condition~3, so their
union is a $\delta$-clique. Conversely suppose that union
is a clique. Take $p\in\ell(x)\setminus\ell(K)$ and
$q\in\ell(y)\setminus\ell(K)$, with carriers $a,b$ in
$x,y$. Neither carrier lies in $K$, so $a,b$ are distinct
exterior strata. Their $\delta$ relation says $p\parallel q$.
Condition~4 yields the original compatibility of $x,y$.

Compatibility of arbitrary figures is determined on their
strata, so it too is recovered. The original deployments
form a permitted chosen atlas: prime generation gives all
downset subfigures, and a clique union of two admitted figures
has pairwise compatible strata and hence is their original
admitted join. Both constructions thus recover the original
orders and figures. They recover the graph on lieux by its
singleton deployments and the full relation by local detection.
All formulas commute with isomorphisms, proving the stated
equivalence, including the one-figure atelier on the empty
stratum set.

For the final assertion, restricted internal transitivity
already gives $U_x\subseteq U_y$ whenever $x\rho y$.
This makes the induced relation hereditary; nonempty open
lieu sets give irreflexivity, and the stipulated graph test
gives principal disjointness. The carrier presentation is
therefore strong. Theorem~\ref{thm:strong-relations} forces
full internal transitivity. Conversely full transitivity is
necessary for every strong system, so the iff follows.
In the factorization presentation transitivity is already
part of the partial-order hypothesis; the same three checks
give the final assertion.

\subsection{General extension and the limits of uniqueness}
We prove Theorem~\ref{thm:abstract-extension}.
If $H\parallel F$ and $P\ll F$, every stratum $p$ of $P$
has a carrier $c$ in $F$. For $h\in H$, disjointness gives
$h\delta c$, and heredity along $p\rho c$ gives
$h\delta p$. Hence $H\parallel P$, so
\[
 \mathcal N(F)\subseteq\mathcal N(P).
\]

Suppose $P\preceq F$ and $P\subsetneq Q\ll F$.
Choose $x\in Q\setminus P$ and a lieu $p\rho x$.
Every stratum $y\in P$ is distinct from $x$ in the clique $Q$,
so $x\delta y$ and then $p\delta y$. Thus $p\parallel P$.
But $p\ll F$, and a nonzero refining lieu cannot be disjoint
from its target: compatible refinement would imply $p\leq F$,
contradicting zero intersection. This contradicts equal cosupports
and proves incidence maximality.

Conversely, if a maximal $P\ll F$ has different cosupport,
witness absorption gives a nonzero $H\ll F$ with
$H\parallel P$ but $H\not\parallel F$.
Its union with $P$ is a proper incidence enlargement refining
$F$, a contradiction. This proves the equivalence.

Now fix $G\leq F$ and $G'\preceq G$.
Order the figures $Q\ll F$ containing $G'$ by inclusion.
This set is nonempty. A chain has a union which is a downset
clique, hence a figure in the full atlas. Every stratum of the
union still refines $F$, so the union is an upper bound in the
same set. Zorn's lemma gives a maximal member $F'$.
Any larger refinement of $F$ would still contain $G'$, so $F'$
is maximal among all refinements of $F$ and hence a subdivision.
Carrier restriction gives
\[
 G'\leq F'|G\ll G.
\]
Maximality of the subdivision $G'$ forces $F'|G=G'$.
This argument also includes the empty figure.

\begin{proposition}[What the deployment square determines]
\label{prop:deployment-uniqueness}
Let $G$ be a downset in a poset $F$, and let $r:G'\to G$
be monotone. An ordered deployment with inverse image $G'$,
the prescribed isomorphism over $F\setminus G$, and the mixed
incidence rule is unique up to unique isomorphism over those data.
It is the poset
\[
 E=G'\amalg(F\setminus G)
\]
with its two original within-piece orders, no exterior-to-interior
relations, and
\[
 \xi\leq\eta\quad\Longleftrightarrow\quad
 r(\xi)\leq_F\eta
 \qquad(\xi\in G',\ \eta\in F\setminus G).
\]
\end{proposition}
\begin{proof}
The map $c:E\to F$, equal to $r$ on $G'$ and identity on the
exterior, is monotone. Mixed transitivity uses either monotonicity
of $r$ or transitivity in $F$. There is no reverse mixed relation,
so antisymmetry holds. The inverse image of $G$ is exactly $G'$.
Every ordered pair in a proposed deployment is prescribed by the
two within-piece orders or the displayed rule. Consequently the
map fixed by the two identifications is the unique isomorphism.
This proves a statement about ordered deployments.
\end{proof}

\begin{counterexample}[Distinct objects with the same extension diagram]
\label{ex:extension-nonunique}
Even a full strong local atlas with witness absorption and all
compatible-cover gluing can have two different elementary
subdivision extensions realizing the same prescribed deployment
square and the same empty boundary.
\end{counterexample}
\begin{proof}
Take strata $a,b$, identity incidence, refinement $a\ll b$,
and empty primitive disjunction. Every principal deployment is a
singleton, carriers are unique, and internal refinement is the
same two-element partial order. Heredity and principal
disjointness are vacuous. The full figures are $0,a,b$.
The unique lieu is $a$, and $U_a=U_b=\{a\}$, so local
detection holds. The two nonzero figures are incompatible and
both have cosupport $\{0\}$.

Witness absorption is vacuous for nonzero source refinements,
because their cosupports agree. For $P=0$ and $F\ne0$,
the witness $F$ itself works. A compatible incidence cover of a
nonzero figure contains that whole figure, so matching restrictions
uniquely determine the glued refinement. This proves gluing for
arbitrary covers, and the empty case is immediate.

For $F=b$ and $G=G'=0$, both $F'=a$ and $F'=b$ are
subdivisions with the required restriction. Their deployments
are one-point posets and their carrier maps are the unique maps
to the one-point deployment of $b$. Both diagrams have the
prescribed empty inverse image, isomorphic complement and vacuous
mixed rule. Both elementary boundaries are empty. The objects
are nevertheless different.

The strengthened diagram extension also exists for every input
in this model. For $(F,G,G')$ equal to
\[
 (0,0,0),\ (a,0,0),\ (a,a,a),\ (b,0,0),\ (b,b,a),\ (b,b,b),
\]
the extension choices are respectively
$0,a,a,\{a\text{ or }b\},a,b$.
These exhaust the possible inputs. The example leaves open a
separate requirement of a coarsest extension; it disproves
uniqueness merely from the deployment square.
\end{proof}
